\begin{problem}[The sheaf of locally polynomial functions]
Let $k$ be a field and $X=k$ with a topology in which polynomial functions make sense (for instance $k=\mathbb R$ with the Euclidean topology). Define $\mathcal P(U)$ to be restrictions to $U$ of a single polynomial in $k[t]$. Describe $\mathcal P^\#$ intrinsically.
\end{problem}

\paragraph{Why this problem matters.}
The example separates global presentation by one formula from local presentation by possibly different formulas.

\paragraph{Useful ideas before starting.}
Use the universal property, stalk preservation, and the fact that a morphism of sheaves is an isomorphism when it is so on every stalk.  Keep the distinction between a presheaf statement on sections and a local statement on germs explicit.

\paragraph{Guided hints.}
Use the locally representable-sections description of sheafification.

\ifdefined\FullProblemDossiers
\begin{solution}
A section of $\mathcal P^\#(U)$ is, by definition of the associated sheaf, a family of germs locally represented by sections of $\mathcal P$. Therefore it is exactly a function $f:U\to k$ such that every point $x\in U$ has a neighborhood $V$ and a polynomial $p_V\in k[t]$ with $f|_V=p_V|_V$. Thus
\[
\mathcal P^\#(U)=\{\text{locally polynomial functions on }U\}.
\]
The canonical map sends one global polynomial formula to the same function viewed locally. On disconnected opens, sheafification permits different polynomial formulas on different components; on connected real intervals, polynomial identity forces compatible local formulas to propagate to one polynomial.
\end{solution}

\paragraph{What happened mathematically.}
Sheafification converts ``one global expression into ``expressible locally in the prescribed class.

\paragraph{Alternative approaches.}
Whenever possible one may replace the espace-etale model by the two-plus construction, or conversely replace a cover/refinement argument by the universal property.  The two approaches agree because both construct the same associated sheaf.

\paragraph{Extensions.}
Formulate the analogous construction for rational functions regular on neighborhoods.

\paragraph{Related problems.}
Compare with VI.14.P1 for the complete legacy construction and with the later sheaf chapters for operations that are deliberately not migrated here.

\paragraph{Provenance.}
New synthesis problem based on the local-representability principle.
\else
\paragraph{Provenance.}
New synthesis problem based on the local-representability principle.
\fi
